\documentclass[aspectratio=169]{beamer}
\input{header}
\coursecode{JKSSB}

\title{Printers \& Displays}
\subtitle{Lesson 15 --- Mechanisms and Display Terminology}
\author{JKSSB Computer Awareness}
\date{\today}

\begin{document}
\frame{\titlepage}

\begin{frame}{Printers: Hard-Copy Output}
  \begin{itemize}
    \item A \key{printer} is an output device that produces \key{hard copy}
      --- output on a physical medium, usually paper.
    \item Output shown on a monitor is called \key{soft copy}; a printer turns
      it into permanent hard copy.
    \item Printers are compared in three main ways: by \key{mechanism}
      (dot matrix, inkjet, laser, thermal, 3D), by \key{impact or non-impact},
      and by the \key{unit printed} (character, line, page).
    \item \muted{Choosing the right labels is the whole exam skill here.}
  \end{itemize}
  \begin{block}{Definition}
    A printer accepts output data from the computer and records it as a
    permanent, human-readable copy on paper or another medium.
  \end{block}
\end{frame}

\begin{frame}{Classifying Printers}
  \begin{center}
    \begin{tabular}{p{0.24\textwidth}p{0.66\textwidth}}
      \toprule
      \textbf{Basis} & \textbf{Types} \\
      \midrule
      By mechanism & Dot matrix, inkjet, laser, thermal, 3D printer \\
      By contact & Impact (strikes the paper) vs non-impact (no contact) \\
      By unit printed & Character printer, line printer, page printer \\
      \bottomrule
    \end{tabular}
  \end{center}
  \vspace{0.25cm}
  \muted{The same printer carries several labels: a laser printer is
  non-impact, a page printer, and uses a laser mechanism.}
\end{frame}

\begin{frame}{Impact vs Non-Impact}
  \begin{columns}[T]
    \column{0.46\textwidth}
      \textbf{Impact printer}
      \begin{itemize}
        \item Forms characters by \key{striking} an inked ribbon against
          the paper
        \item Noisy; can print carbon copies
        \item Dot matrix, daisy wheel, drum/line printer
      \end{itemize}
    \column{0.46\textwidth}
      \textbf{Non-impact printer}
      \begin{itemize}
        \item Prints \key{without striking} the paper
        \item Quieter; better quality; no carbon copies
        \item Inkjet, laser, thermal
      \end{itemize}
  \end{columns}
  \vspace{0.3cm}
  \begin{alertblock}{Trap alert}
    Impact is about \key{striking the paper}, not about how the image is
    formed. Dot matrix strikes the paper; inkjet, laser, and thermal do not.
  \end{alertblock}
\end{frame}

\begin{frame}{Impact Printers}
  \begin{itemize}
    \item An impact printer forms a character by \key{physically hitting} an
      inked ribbon against the paper.
    \item \key{Dot matrix}: pins hit the ribbon to build characters from dots.
    \item \key{Daisy wheel}: a spinning wheel of type characters, like a
      typewriter; good text, no graphics.
    \item \key{Line printers} such as the drum printer and chain printer
      strike a whole line at a time at high speed.
    \item Because it presses through, an impact printer can make
      \key{multiple carbon copies} of multi-part forms.
  \end{itemize}
\end{frame}

\begin{frame}{Dot Matrix Printer}
  \begin{itemize}
    \item It is an \key{impact} printer. The \key{print head} contains rows
      of tiny pins.
    \item The pins strike an \key{inked ribbon} to form each character as a
      pattern of dots in a matrix.
    \item It is \key{noisy}, cheap to run, and can print \key{carbon copies}
      and multi-part invoices.
    \item Print quality is low and graphics are poor; still used for receipts
      and billing.
  \end{itemize}
  \begin{alertblock}{Trap alert}
    A dot matrix printer is an \key{impact} printer (it strikes the ribbon).
    Do not group it with inkjet, laser, or thermal.
  \end{alertblock}
\end{frame}

\begin{frame}{Line Printers}
  \begin{itemize}
    \item A \key{line printer} prints an entire \key{line} of text in one
      operation, not one character at a time.
    \item \key{Drum printer}: a rotating drum carries the character set; a
      hammer strikes the paper where each character is needed.
    \item \key{Chain (band) printer}: a moving chain or band of type
      characters runs past the hammers.
    \item They are high-speed \key{impact} printers used with older
      mainframes and large batch jobs.
  \end{itemize}
  \begin{block}{Print units}
    Character printer $\rightarrow$ line printer $\rightarrow$ page printer,
    from smallest to largest unit printed.
  \end{block}
\end{frame}

\begin{frame}{Inkjet Printer}
  \begin{itemize}
    \item A \key{non-impact} printer: it does not touch the paper.
    \item The print head sprays tiny \key{droplets of liquid ink} through
      fine nozzles onto the page.
    \item It is quieter than an impact printer and prints good-quality
      \key{full colour} images on many paper types.
    \item It is slower per page than a laser printer, and wet ink can
      \key{smudge}; ink cartridges need replacing.
  \end{itemize}
  \begin{block}{Remember}
    Inkjet = \key{liquid ink} sprayed through nozzles; no toner, no striking.
  \end{block}
\end{frame}

\begin{frame}{Laser Printer}
  \begin{itemize}
    \item A \key{non-impact} printer that uses a \key{laser beam} and
      \key{toner} (fine powder) in a xerographic process.
    \item The laser draws an image on a charged drum; toner sticks to it and
      is then fused onto the paper with heat and pressure.
    \item It is \key{fast} with sharp, high-quality text, and it is a
      \key{page printer} --- it prints a whole page at a time.
    \item Toner is used instead of ink, so output does not smudge when wet.
  \end{itemize}
  \begin{alertblock}{Trap alert}
    Laser = \key{toner} + laser beam. Inkjet = \key{liquid ink}. Do not swap
    the two consumables.
  \end{alertblock}
\end{frame}

\begin{frame}{Thermal Printer}
  \begin{itemize}
    \item A \key{non-impact} printer that produces an image using \key{heat}.
    \item It heats specially treated \key{thermal paper} (or a heat-sensitive
      ribbon) to form characters.
    \item It is \key{silent} in operation and needs no ink or toner
      cartridge in the simplest design.
    \item Used for POS receipts, tickets, barcode labels, and older fax
      machines; usually limited to one colour.
  \end{itemize}
  \begin{block}{Keep it non-impact}
    Thermal, laser, and inkjet are \key{non-impact}. Dot matrix is
    \key{impact}.
  \end{block}
\end{frame}

\begin{frame}{3D Printer}
  \begin{itemize}
    \item A \key{3D printer} builds a solid three-dimensional object
      \key{layer by layer} from a digital model --- this is
      \key{additive manufacturing}.
    \item It does not print ink on paper; it adds material (plastic, resin,
      or metal) until the object is formed.
    \item Used for prototypes, models, spare parts, and medical implants.
    \item \muted{The word ``printer'' is used, but the output is a physical
      object, not a document.}
  \end{itemize}
\end{frame}

\begin{frame}{Plotters}
  \begin{itemize}
    \item A \key{plotter} is an output device that \key{draws} line
      graphics, usually with one or more pens moving over the paper.
    \item It works with \key{vector} drawings, so it is ideal for large
      \key{engineering drawings}, maps, and CAD plans.
    \item Plotters handle big sheet sizes that ordinary printers cannot.
    \item \muted{A plotter draws lines; a printer lays down dots or
      characters. Both are output devices.}
  \end{itemize}
  \begin{alertblock}{Trap alert}
    A plotter is an \key{output} device, not an input device. Do not confuse
    it with a scanner or a digitizer (TRAP-12).
  \end{alertblock}
\end{frame}

\begin{frame}{Inkjet vs Laser (Near-Neighbour)}
  \begin{center}
    \begin{tabular}{lll}
      \toprule
      \textbf{Point} & \textbf{Inkjet} & \textbf{Laser} \\
      \midrule
      Contact & Non-impact & Non-impact \\
      Image formed by & Liquid ink droplets & Laser beam + toner \\
      Unit printed & Character/line & Page printer \\
      Speed & Slower & Faster \\
      Strengths & Colour photos & Sharp text, offices \\
      \bottomrule
    \end{tabular}
  \end{center}
  \vspace{0.25cm}
  \muted{Both are non-impact. The usual trap swaps their consumables and
  speeds.}
\end{frame}

\begin{frame}{Impact or Non-Impact? Sort Each Printer}
  \begin{itemize}
    \item \key{Impact}: dot matrix, daisy wheel, drum printer, chain printer.
    \item \key{Non-impact}: inkjet, laser, thermal.
    \item A \key{page printer} (such as a laser printer) prints a whole page
      at a time; a \key{line printer} prints a whole line at a time.
    \item \muted{3D printers and plotters are in their own categories; they
      are not ordinary paper printers.}
  \end{itemize}
  \begin{alertblock}{Trap alert}
    Only devices that \key{strike an inked ribbon} are impact. Inkjet,
    laser, and thermal never touch the paper.
  \end{alertblock}
\end{frame}

\begin{frame}{Display Terminology: Pixel and Resolution}
  \begin{itemize}
    \item A \key{pixel} (picture element) is the \key{smallest addressable
      element} of a display; every image is built from pixels.
    \item \key{Resolution} is the number of pixels the display can show,
      written as columns $\times$ rows (for example $1920 \times 1080$).
    \item \key{Higher resolution} means more pixels and a \key{sharper}
      image.
    \item Printer sharpness is measured in \key{DPI} =
      \key{Dots Per Inch}.
  \end{itemize}
  \begin{block}{One idea}
    Pixels are the units; resolution is how many of them fill the screen.
  \end{block}
\end{frame}

\begin{frame}{Display Terminology: Refresh Rate and Response Time}
  \begin{itemize}
    \item \key{Refresh rate} is how many times per second the display
      redraws the image; measured in \key{hertz (Hz)}.
    \item A \key{higher refresh rate} (for example 120 Hz) gives smoother
      motion and less flicker.
    \item \key{Response time} is how long a pixel takes to change from one
      state to another; measured in \key{milliseconds (ms)}.
    \item A \key{lower response time} means less motion blur or ghosting.
  \end{itemize}
  \begin{alertblock}{Trap alert}
    Refresh rate is in \key{Hz} (higher is better); response time is in
    \key{ms} (lower is better). Do not swap the units.
  \end{alertblock}
\end{frame}

\begin{frame}{Brightness and Contrast}
  \begin{itemize}
    \item \key{Brightness} is the amount of light a display emits; it is
      commonly measured in \key{nits} (candela per square metre, cd/m$^2$).
    \item \key{Contrast} is the difference between the brightest white and
      the darkest black a display can produce.
    \item The \key{contrast ratio} compares the two (for example 1000:1);
      a higher ratio gives deeper blacks and more vivid images.
    \item \muted{Brightness is how much light; contrast is the range between
      light and dark.}
  \end{itemize}
  \begin{block}{Summary of units}
    Resolution $\rightarrow$ pixels/DPI \;$\cdot$\; refresh rate $\rightarrow$
    Hz \;$\cdot$\; response time $\rightarrow$ ms \;$\cdot$\;
    brightness $\rightarrow$ nits.
  \end{block}
\end{frame}

\begin{frame}{Lesson 15: Recall}
  \begin{itemize}
    \item A printer produces \key{hard copy}; monitor output is soft copy.
    \item \key{Impact} printers (dot matrix, daisy wheel, drum/chain line
      printers) strike an inked ribbon; \key{non-impact} printers (inkjet,
      laser, thermal) do not.
    \item Dot matrix = impact, pins + ribbon, carbon copies; inkjet =
      liquid ink sprayed; laser = laser beam + toner, a page printer;
      thermal = heat-sensitive paper.
    \item A 3D printer adds material layer by layer; a plotter draws vector
      line graphics for engineering drawings.
    \item Display terms: \key{pixel} (smallest element), \key{resolution}
      (pixels shown), \key{refresh rate} (Hz), \key{response time} (ms),
      \key{brightness} (nits), \key{contrast} (bright/dark ratio).
  \end{itemize}
\end{frame}
\end{document}
